\section{Preliminaries}
\label{sec:preliminaries}

This section assembles the minimal background needed to read the rest of the
paper. The exposition deliberately addresses two audiences in parallel:
readers fluent in topological data analysis (TDA) who have only passing
familiarity with formal causal inference (CI), and readers fluent in CI who
have only passing familiarity with persistent homology. We therefore strive
for completeness over compression and fix the notation that
\Cref{sec:framework,sec:literature,sec:case_studies,sec:research} will reuse.

\subsection{Topological data analysis}
\label{sec:prelim_tda}

TDA studies the multi-scale geometric and topological structure of data
using tools from algebraic topology
\citep{carlsson2009topology,edelsbrunner2010computational,
chazal2021introduction,wasserman2018topological}. The unifying idea is
\emph{functoriality}: a data set is mapped, in a stable and computable way,
to an algebraic object whose invariants summarise its shape.

\subsubsection{Simplicial complexes and filtrations}

Let $V$ be a finite vertex set. An abstract simplicial complex on $V$ is a
collection $K \subseteq 2^V$ closed under taking subsets: if $\sigma \in K$
and $\tau \subseteq \sigma$, then $\tau \in K$. A simplex
$\sigma \in K$ with $|\sigma| = k+1$ is called a $k$-simplex; the dimension
of $K$ is the maximum dimension of its simplices. A 1-dimensional simplicial
complex is a graph; higher-dimensional simplicial complexes admit triangles,
tetrahedra, and so on.

Given a finite point cloud $X \subset (\Xc, d)$ in a metric space and a
scale parameter $r \ge 0$, the \emph{Vietoris--Rips complex} is
\begin{equation}
  \mathrm{Rips}_r(X)
  = \bigl\{ \sigma \subseteq X : d(u,u') \le 2r \text{ for all } u, u' \in \sigma \bigr\},
  \label{eq:rips}
\end{equation}
and the \emph{\v{C}ech complex} $\mathrm{\check{C}}_r(X)$ collects subsets
whose closed $r$-balls have non-empty common intersection
\citep{hausmann2016vietoris}. For $X \subset \R^m$ the \emph{$\alpha$-complex}
$\mathrm{Alpha}_r(X)$ is built from intersections of Voronoi cells with
$r$-balls and yields a homotopy-equivalent but smaller complex than
$\mathrm{\check{C}}_r$. For image-like grid data, \emph{cubical complexes}
play the analogous role, with elementary cubes replacing simplices.

A \emph{filtration} of $K$ is a nested family
$\Filt = \{K(t)\}_{t \in \R}$ with $K(s) \subseteq K(t)$ whenever
$s \le t$. Increasing $r$ in \cref{eq:rips} produces a Vietoris--Rips
filtration; sublevel sets $K(t) = f^{-1}((-\infty,t])$ of a monotone
function $f : K \to \R$ produce a sublevel filtration. Filtrations
expose the data's structure at multiple resolutions simultaneously.

\subsubsection{Homology and persistent homology}

Singular or simplicial homology assigns to each topological space (or
simplicial complex) $X$ a graded sequence of abelian groups
$H_k(X)$, where $H_0$ counts connected components, $H_1$ counts independent
1-dimensional loops, $H_2$ counts enclosed voids, and so on; the rank of
$H_k(X)$ is the $k$th \emph{Betti number} $\betti_k(X)$
\citep{hatcher2002algebraic,munkres1984elements}. Homology is a functor: a
continuous map $f : X \to Y$ induces a group homomorphism
$H_k(f) : H_k(X) \to H_k(Y)$.

Applied to a filtration $\Filt = \{K(t)\}$, the inclusion
$K(s) \hookrightarrow K(t)$ induces, for every $k$, a linear map
$H_k(K(s)) \to H_k(K(t))$. The collection
$\PHk{k}(\Filt) = \{H_k(K(t))\}_{t \in \R}$ together with these maps is the
\emph{$k$-th persistent homology module} of $\Filt$
\citep{edelsbrunner2002topological,zomorodian2005computing,
edelsbrunner2008persistent}. Over a field, any tame persistent homology
module decomposes uniquely (up to isomorphism) into a direct sum of interval
modules, and the corresponding intervals $[b,d) \subset \R$ form the
\emph{persistence barcode}; equivalently one may record their endpoints as
a multiset of points
\begin{equation}
  \Dgm_k(\Filt) \subset \R^2_+ := \{ (b,d) \in (\R \cup \{\infty\})^2 : b < d \},
  \label{eq:dgm}
\end{equation}
the \emph{persistence diagram} in homology degree $k$
\citep{ghrist2008barcodes,edelsbrunner2010computational}. Following standard
convention, $\Dgm_k$ includes the diagonal $\Delta = \{(t,t) : t \in \R\}$
with infinite multiplicity, so matchings between diagrams may send unmatched
points to $\Delta$.

\subsubsection{Distances on diagrams and the stability theorem}

The two principal metrics on the space of persistence diagrams are the
\emph{bottleneck distance} $\dB$ and the $p$-Wasserstein distance $W_p$:
\begin{equation}
  \dB(D, D') = \inf_{\gamma} \, \sup_{x \in D} \|x - \gamma(x)\|_\infty,
  \qquad
  W_p(D, D') = \Bigl( \inf_{\gamma} \sum_{x \in D} \|x - \gamma(x)\|_\infty^p \Bigr)^{1/p},
  \label{eq:bottleneck}
\end{equation}
where $\gamma$ ranges over bijections $D \cup \Delta \to D' \cup \Delta$.
The cornerstone of statistical TDA is the \emph{stability theorem}: if $f,g$
are tame functions on a triangulable space, then
\begin{equation}
  \dB\bigl(\Dgm_k(f), \Dgm_k(g)\bigr) \le \|f - g\|_\infty,
  \label{eq:stability}
\end{equation}
and analogous bounds hold in Wasserstein distance
\citep{cohen2007stability,chazal2016structure}. Stability is the property
that makes persistent homology a viable summary for noisy data and is what
will allow us, in \Cref{sec:framework}, to lift TDA-level perturbation
guarantees to causal-inference-level robustness statements.

\subsubsection{Functional summaries: landscapes, silhouettes, Euler curves}

For statistical and machine-learning use, it is convenient to embed
$\Dgm_k$ in a Hilbert or Banach space. Three standard embeddings will
appear repeatedly. For $p = (a,b) \in \Dgm_k$, define the tent function
$\Lambda_p(t) = \max\{0, \min(t-a, b-t)\}$. The
\emph{persistence landscape}
$\{\lambda_j(t; D)\}_{j \ge 1}$ records the $j$-th largest value of
$\{\Lambda_p(t)\}_{p \in D}$ \citep{bubenik2015statistical}. The
\emph{power-weighted silhouette} of order $r > 0$ is
\begin{equation}
  \varphi(t; D, r)
  = \frac{\sum_{p = (a,b) \in D} (b - a)^r \, \Lambda_p(t)}
         {\sum_{p = (a,b) \in D} (b - a)^r}, \qquad t \in T \subset \R,
  \label{eq:silhouette}
\end{equation}
with $T$ a fixed compact interval
\citep{chazal2014stochastic,kim2026topological}. The \emph{persistence
image} \citep{adams2017persistence} provides a fixed-dimensional vector
embedding. Finally, the \emph{Euler characteristic curve} records
$\chi(K(t)) = \sum_k (-1)^k \betti_k(K(t))$ along the filtration; though it
discards information, it is cheap to compute and frequently appears in
applications \citep{smith2025multiscale,smith2025harnessing}.

\begin{lemma}[Silhouette Lipschitz stability; informal restatement of
\citealp{kim2026topological}]
\label{lem:silhouette_stability}
Under a uniform bound on tent heights, the power-weighted silhouette
satisfies
$\|\varphi(\cdot; D) - \varphi(\cdot; D')\|_\infty \le c_r \, W_1(D, D')$
for an explicit constant $c_r > 0$ depending only on $r$ and an upper bound
on persistence lifetimes.
\end{lemma}

\Cref{lem:silhouette_stability} is the bridge by which Wasserstein
perturbations of diagrams translate into sup-norm perturbations of
silhouettes; \Cref{sec:framework} uses it to derive causal-stability
statements.

\subsubsection{Mapper and Reeb graphs}

The Mapper algorithm
\citep{singh2007mapper,madukpe2026mapper,dlotko2019ball} produces a
combinatorial summary of a high-dimensional point cloud $X$. Given a
\emph{filter} (or lens) function $f : X \to \R^d$, a cover
$\{U_\alpha\}$ of $f(X)$, and a clustering procedure, Mapper assigns one
node per cluster of $f^{-1}(U_\alpha)$ and connects nodes whose clusters
intersect. The resulting simplicial complex is a discrete proxy for the
Reeb graph of $f$, which collapses connected components of level sets to
points
\citep{biasotti2008reeb,bauer2013measuring,morozov2013interleaving}. Reeb
graphs and Mapper are central in the topological abstraction of causal
mechanisms (\Cref{sec:framework}, Domain A2).

\subsection{Causal inference}
\label{sec:prelim_ci}

Causal inference seeks to quantify what would happen to an outcome under
hypothetical interventions on a system, going beyond the associations
captured by joint distributions
\citep{pearl2009causality,imbens2015causal,peters2017elements,
spirtes2001causation}. Two complementary frameworks dominate the
methodological landscape.

\subsubsection{Structural causal models}

A \emph{structural causal model} (SCM) is a tuple
$\Mc = \langle U, V, \{f_V\}_{V \in V}, P \rangle$, where $V$ is a set of
\emph{endogenous} variables, $U$ a set of \emph{exogenous} variables,
$f_V : \chi_{\Pa(V)} \times \chi_{U(V)} \to \chi_V$ is a measurable
structural function determining each endogenous variable from its parents
and exogenous noise, and $P$ is a Borel probability measure on $\chi_U$
\citep{pearl2009causality,bongers2021foundations,
ibeling2021topological}. Each SCM induces a unique observational
distribution $p^\Mc \in \Prob(\chi_V)$ under recursiveness assumptions;
under an \emph{intervention} $W := w$ on a subset $W \subseteq V$, the
intervened SCM $\Mc_{W := w}$ holds $W = w$ fixed and otherwise leaves
mechanisms intact, yielding an \emph{interventional} distribution
$p^{\Mc_{W := w}}$ \citep{pearl2009causality}.

\emph{Pearl's do-calculus} \citep{pearl1995causal,tian2002general,
lee2019general} provides a complete set of rewrite rules for the
intervention operator $\doop(\cdot)$ on a directed acyclic graph (DAG); it
permits the algorithmic identification of interventional from
observational distributions whenever such identification is possible. The
\emph{causal hierarchy}---observational, interventional, counterfactual---
organises causal queries by the strength of assumption needed to identify
them \citep{bareinboim2020pearls,ibeling2021topological}.

A \emph{causal graph} $G = (V, E)$ on the variable set $V$ encodes the
direct-influence relation: $X \to Y \in E$ means $X \in \Pa(Y)$. We write
$G$ for a generic causal graph and reserve $\Mc$ for SCMs whose graph
structure is $G$. The space of all SCMs compatible with $G$ is denoted
$\mathfrak{M}_G$; without graph restriction we write $\mathfrak{M}$. The
topology on $\mathfrak{M}$ used in \citet{ibeling2021topological} is the
weak topology induced by the convergence of the joint distribution over
all interventions, and is one of the central objects in
\Cref{sec:framework}, Domain B.

\subsubsection{Potential outcomes and causal estimands}

The \emph{potential outcomes} (or Neyman--Rubin) framework
\citep{rubin1974estimating,rubin2005causal,imbens2015causal} parameterises
causal effects directly in terms of counterfactual outcomes. For a binary
treatment $A \in \{0,1\}$ and observed outcome $Y$, one posits
counterfactual outcomes $Y^0, Y^1$ representing the response under each
treatment. The \emph{average treatment effect} is
\begin{equation}
  \ATE = \E[Y^1 - Y^0],
  \label{eq:ate}
\end{equation}
and conditioning on covariates $X$ yields the \emph{conditional average
treatment effect} $\CATE(x) = \E[Y^1 - Y^0 \mid X = x]$. Identification of
\cref{eq:ate} from observational data typically rests on three
assumptions:
\begin{enumerate}[label=(C\arabic*), leftmargin=*]
  \item \emph{Consistency}: $Y = Y^A$ almost surely.
  \item \emph{No unmeasured confounding} (a.k.a.\ ignorability):
        $\{Y^0, Y^1\} \ind A \mid X$.
  \item \emph{Positivity (overlap)}: $0 < \Prob(A = 1 \mid X) < 1$ almost surely.
\end{enumerate}
Under (C1)--(C3), the ATE is identified by
$\ATE = \E[\mu_1(X) - \mu_0(X)]$ with $\mu_a(x) = \E[Y \mid X = x, A = a]$,
and is estimable by inverse-probability weighting, outcome regression, or
\emph{augmented inverse-probability weighting} (AIPW)
\citep{kennedy2024semiparametric,chernozhukov2018double}. AIPW is
\emph{doubly robust}: consistency holds if either the outcome or
propensity model is consistent, and rate-double-robust efficiency holds
under product-rate nuisance conditions
\citep{kennedy2023towards,newey2018cross}.

When the outcome is functional or non-Euclidean, the ATE in \cref{eq:ate}
is no longer well defined, and one must replace the difference of means by
a metric- or geometry-respecting estimand. Two such generalisations recur
in this paper:
\begin{itemize}[leftmargin=*]
  \item the \emph{functional average treatment effect}, defined
        pointwise in a function space and estimated via doubly-robust
        AIPW on the function-on-scalar regression
        \citep{ecker2024causal,testa2025doubly,liebl2023fast};
  \item the \emph{geodesic average treatment effect}, defined in a
        geodesic metric space and estimated via Fréchet means
        \citep{kurisu2024geodesic}.
\end{itemize}
Functional ATEs on \emph{topological} summaries (such as silhouettes) are
exactly what \citet{kim2026topological} formalise, and we will see in
\Cref{sec:framework} that this construction is one of three founding
pillars of the TCDA framework.

\subsubsection{Causal discovery}

Beyond effect estimation, \emph{causal discovery} attempts to learn the
graph $G$ itself from data
\citep{spirtes2001causation,peters2014causal,
shimizu2006lingam,causaldisco2023survey}. Methods fall into three broad
families:
\begin{itemize}[leftmargin=*]
  \item \emph{Constraint-based}: query conditional independencies and
        prune candidate DAGs accordingly (PC, FCI).
  \item \emph{Score-based}: define a regularised likelihood over DAGs and
        search the space (GES, NOTEARS).
  \item \emph{Functional causal models}: exploit functional asymmetries
        such as additive non-Gaussian noise to identify the causal
        direction
        (LiNGAM, ANM).
\end{itemize}
A growing body of work uses \emph{topological orderings} of variables to
shrink the search space
\citep{xu2025information,tran2024constraining,wu2026optimization,
barroso2024inference}; this line is one of the most direct existing
bridges between TDA and CI and is taken up in \Cref{sec:framework},
Domain A1.

\subsection{Existing mathematical bridges}
\label{sec:prelim_bridges}

The interface between TDA and CI is not a void: several mathematical
threads already connect the two communities, although they have not yet
been unified into a single domain. We review three.

\subsubsection{Topologies on causal-model spaces}

Following \citet{dembo1994topological} and \citet{genin2017topology},
open sets of the weak topology on probability distributions correspond to
statistically \emph{verifiable} hypotheses. \citet{ibeling2021topological}
extend this correspondence to the space of SCMs by constructing a sequence
of topologies, one for each level of the causal hierarchy, and proving
that the collapse set of the hierarchy is topologically meager. Earlier
parametric results on faithfulness violations
\citep{meek1995strong,uhler2013geometry} can be re-read in this light as
measure-theoretic predecessors to the topological statement. We use the
phrase ``topology of causal models'' for this line of work and adopt it
as the foundation of \Cref{sec:framework}, Domain B.

\subsubsection{Sheaf-theoretic and categorical approaches}

Sheaves, originally developed for cohomology and complex geometry
\citep{bredon1997sheaf}, have re-emerged as a language for distributed
data, signal processing, and now causal inference
\citep{robinson2014topological,curry2014sheaves,spivak2014category}.
\citet{mahadevan2021causal,mahadevan2022universal,mahadevan2025sheaf}
develop a topos-theoretic formulation of causal inference in which
do-calculus is generalised via subobject classifiers, the space of causal
DAGs is endowed with an Alexandroff topology, and certain identification
results take the form of sheaf-cohomology statements. Independently, the
categorical view of persistent homology---homology is a functor, and
filtrations are functors from $(\R, \le)$ to the category of topological
spaces---suggests that any \emph{functorial} causal inference procedure
should commute, up to a controllable error, with the topological functor
used to summarise data. This is the formal content of Domain B3
(``functorial causal inference'') in \Cref{sec:framework}.

\subsubsection{Optimal transport and Wasserstein space}

Optimal transport equips spaces of probability measures with rich
geometric structure
\citep{villani2009optimal,peyre2019computational}. Two threads concern
this paper. First, persistence diagrams are themselves compared via
$W_p$-distances; many statistical guarantees for persistent homology are
expressed in these terms
\citep{mileyko2011probability,turner2014frechet,
chazal2014stochastic}. Second, the Wasserstein distance offers a natural
metric on counterfactual distributions and so on causal effect
estimators, with stability results connecting changes in observed
distributions to changes in identified causal quantities. The combination
of these two threads underwrites the silhouette stability bound of
\citet{kim2026topological}, which we will use in
\Cref{sec:framework}.

\subsection{Notation summary}
\label{sec:notation}

We collect the notation that recurs throughout the paper.

\begin{center}
\begin{tabular}{@{}lp{0.78\linewidth}@{}}
\toprule
Symbol & Meaning \\
\midrule
$(\Mc, d)$       & Topological / metric space of observations (a Polish space, in the broadest setting). \\
$\Filt = \{K(t)\}_{t \in \R}$ & A filtration of simplicial / cubical complexes built from data. \\
$\PHk{k}(\Filt)$ & $k$-th persistent homology module of the filtration. \\
$\Dgm_k(\Filt)$  & $k$-th persistence diagram (multiset of birth--death pairs). \\
$\dB, \dW, W_p$  & Bottleneck and $p$-Wasserstein distances on diagrams. \\
$\varphi(t; D, r)$ & Power-weighted silhouette function of diagram $D$ at exponent $r$. \\
$\Map_f, \Reeb_f$ & Mapper and Reeb-graph summaries with filter $f$. \\
$G = (V, E)$     & Causal graph (directed acyclic, unless stated). \\
$\Mc \in \mathfrak{M}_G$ & SCM with graph $G$; $\mathfrak{M}$ if the graph is unspecified. \\
$p^\Mc, p^{\Mc_{W := w}}$ & Observational and interventional distributions of $\Mc$. \\
$Y^a$             & Potential outcome under treatment $A = a$. \\
$\ATE, \CATE, \TATE, \GATE$ & Average / conditional / topological / geodesic ATEs. \\
$\Tc$            & A topological functor (persistent homology, Mapper, etc.) acting on data. \\
$\Cc$            & A causal query (e.g., $\ATE$, $\doop$-expression, identifiability question). \\
$\TCDA$          & Topological Causal Data Analysis (the proposed domain). \\
\bottomrule
\end{tabular}
\end{center}

This notation is used uniformly in \Cref{sec:framework} onwards. The
quadruple $(\Mc, G, \Tc, \Cc)$ in particular will be the carrier of the
TCDA formalism.

\medskip

Equipped with these tools, we now turn to the main intellectual
contribution of the paper.
